% =========================================================
% 4.6 MODELO DE EVALUACIÓN CRIMINOLÓGICA
% El \section{} de esta sección está en el chapter.tex del capítulo;
% sus referencias van en seccion4_6.bib (misma carpeta).
% =========================================================

El presente proyecto estudia únicamente delitos violentos y se fundamenta principalmente en la teoría de patrones delictivos, complementada con otras teorías de criminología ambiental\cite{brantingham1993}, que sustentan la construcción del modelo.

\subsection{Definición de variables}

A partir de la revisión de literatura y de los antecedentes analizados, se definió un conjunto de variables candidatas, organizadas en cuatro categorías: variables físicas del entorno, variables de percepción, variables de control socioeconómico y variables de control por generadores de crimen. Su selección definitiva estará sujeta a la disponibilidad de datos para el área de estudio.

\subsubsection{Variable dependiente}
Determina los datos reales y reportados con los cuales se hará el cálculo de la relevancia por variable. Se utiliza únicamente la variable mostrada en la tabla \ref{tab:variable-dependiente}.

\begin{table}[h]
\centering
\begin{tabular}{|p{3cm}|p{6cm}|p{3cm}|}
\hline
\textbf{Variable} & \textbf{Descripción} & \textbf{Unidad} \\ \hline
Incidencia delictiva & Número de delitos violentos reportados por segmento de calle dentro del periodo de estudio & Conteo \\ \hline
\end{tabular}
\caption{Variable dependiente del modelo}
\label{tab:variable-dependiente}
\end{table}

\subsubsection{Variables físicas del entorno construido}
Funcionan como una serie de variables a nivel micro que definen aspectos físicos de las calles. Estas se muestran en la tabla \ref{tab:variables-fisicas}.

\begin{table}[h]
\centering
\begin{tabular}{|p{3cm}|p{6cm}|p{3cm}|}
\hline
\textbf{Variable} & \textbf{Descripción} & \textbf{Fuente} \\ \hline
Iluminación & Infraestructura de alumbrado del segmento: presencia y densidad de luminarias (escala 1--4) & Street View / crowdsourced \\ \hline
Apertura / visibilidad & Presencia de obstrucciones visuales, muros ciegos o esquinas sin línea de visión & Street View \\ \hline
Vía peatonal & Estado de la banqueta (ninguna, mala, regular, buena) & Street View / crowdsourced \\ \hline
Vegetación & Porcentaje de la imagen ocupado por vegetación & Street View \\ \hline
Presencia de personas & Porcentaje o conteo aproximado de personas visibles & Street View / crowdsourced \\ \hline
Mantenimiento del entorno & Presencia de grafiti, basura o edificios en mal estado & Street View \\ \hline
Uso de suelo & Categoría del uso de suelo (residencial, comercial, mixto) & OSM / catastro \\ \hline
\end{tabular}
\caption{Variables físicas derivadas del entorno construido}
\label{tab:variables-fisicas}
\end{table}

Las variables extraídas de Street View reflejan las condiciones diurnas de un instante determinado, por lo que no capturan variaciones horarias del entorno (por ejemplo, la iluminación efectiva durante la noche o la presencia habitual de personas). Esta limitación se considera en la interpretación de los resultados.

\subsubsection{Variables de percepción}
Recogen cómo los individuos perciben su entorno y cómo lo hace el modelo de inteligencia artificial (\textit{describir el modelo empleado}). Su salida se contrasta con una muestra de segmentos etiquetada manualmente para verificar su fiabilidad. Estas variables se establecen en la tabla \ref{tab:variables-percepcion}.

\begin{table}[h]
\centering
\begin{tabular}{|p{3cm}|p{6cm}|p{3cm}|}
\hline
\textbf{Variable} & \textbf{Descripción} & \textbf{Fuente} \\ \hline
Percepción de seguridad & Grado en que el segmento es percibido como seguro & Modelo de percepción \\ \hline
Percepción de abandono & Grado en que el segmento es percibido como deteriorado o descuidado & Modelo de percepción \\ \hline
\end{tabular}
\caption{Variables de percepción del entorno}
\label{tab:variables-percepcion}
\end{table}

\subsubsection{Variables de control por generadores de crimen}

Con el fin de aislar el efecto de las variables de criminología ambiental de otros factores conocidos como generadores de crimen, se contempla la inclusión de las siguientes variables de control, calculadas mediante un área de influencia alrededor de cada segmento, con decaimiento exponencial en función de la distancia. Estas se muestran en la tabla \ref{tab:variables-control}.

\begin{table}[h]
\centering
\begin{tabular}{|p{3cm}|p{6cm}|p{3cm}|}
\hline
\textbf{Variable} & \textbf{Descripción} & \textbf{Fuente} \\ \hline
Proximidad a expendios de alcohol & Distancia ponderada a bares o tiendas de conveniencia con venta de alcohol & OSM / Google Places API \\ \hline
Proximidad a lotes baldíos & Distancia ponderada a predios sin uso o abandonados & Street View / catastro \\ \hline
Proximidad a transporte público & Distancia ponderada a paradas de metro, metrobús o autobús & OSM / datos públicos \\ \hline
Densidad de actividad comercial & Número de comercios dentro del área de influencia del segmento & OSM / Google Places API \\ \hline
\end{tabular}
\caption{Variables de control por generadores de crimen}
\label{tab:variables-control}
\end{table}

\subsubsection{Variables de control socioeconómico}

Adaptadas al contexto mexicano a partir de datos disponibles a nivel de Área Geoestadística Básica (AGEB) publicados por el INEGI, con el fin de controlar el efecto de la desventaja social concentrada sobre la incidencia delictiva. Estas variables se muestran en la tabla \ref{tab:variables-control-socioeconomico}.

\begin{table}[h]
\centering
\begin{tabular}{|p{3cm}|p{6cm}|p{3cm}|}
\hline
\textbf{Variable} & \textbf{Descripción} & \textbf{Fuente} \\ \hline
Nivel educativo & Porcentaje de población con educación superior en el AGEB & INEGI \\ \hline
Densidad poblacional & Población total por AGEB & INEGI \\ \hline
Grupo etario predominante & Porcentaje de población entre 15 y 29 años & INEGI \\ \hline
\end{tabular}
\caption{Variables de control socioeconómico}
\label{tab:variables-control-socioeconomico}
\end{table}

\subsubsection{Nota metodológica sobre la dirección de las relaciones}

La dirección e intensidad de todas las variables será determinada experimentalmente durante la fase de construcción del modelo de criminología ambiental para el prototipo con base en lo que establece el modelado del terreno de riesgo.

\subsection{Ponderación de los factores}
El método de ponderación de las variables elegidas es totalmente empírico. Se utilizarán los datos históricos de incidencia delictiva real para estimar la relevancia de cada variable y, en algunos casos, descartarla, para así construir un modelo. Los coeficientes resultantes del modelo estadístico determinarán tanto la significancia estadística como el peso específico de cada factor dentro del índice final, garantizando que la ponderación refleje las condiciones empíricas del entorno analizado.
